\section{Introduction}
% Closed-loop Photorealistic driving simulator with editing ability (change the weather or appearance of some objects) is beneficial for training the end-to-end autonomous driving systems. Graphics-based driving simulator carla [XXX]... can achieve full control of the driving scene like the free-view pose rendering, weather editing, arbitraty cars manipulation, etc. However, the rendered results from Graphics-based driving simulator are not photorealistic and have huge gap with the real world driving scenes. Reconstruction based simulator like PVG, OmniRe,....[] can only achieve photo-realistic results within the recorded trajectory, which makes it difficult to be a Closed-loop driving simulator. To solve the off-trajectory problem, some methods freesim, dd4d, recon,[XXX] utilizes the diffusion model to generate novel trajectory views and reconstruct the dynamic scene using the original and generated views. However, the generative model may generate inconsistent things on the same unseen region during different novel trajectories if these inconsistent novel views and original recorded view are used together to reconstruct the driving scene. The results may be blurred. And these reconstruction based simulators can not do the editing tasks. Current world model Cosmos, MaskGWM, etc [XXX] can only generate videos without constructing 4D world. Thus there are not closed-loop simulator and can not control the camera pose explicitly. 



Developing a photorealistic driving simulator with scene editing capabilities is critical for training and evaluating end-to-end autonomous driving systems. 
% Such simulators offer a safe, controllable, and diverse environment for testing perception, planning, and control modules under varying conditions. 
% Traditional graphics-based simulators, such as CARLA~\cite{dosovitskiy2017carla} and Airsim~\cite {shah2018airsim}, provide full control over the scene, including free-viewpoint rendering, weather manipulation, and arbitrary object placement. 
% However, despite their flexibility, these simulators suffer from a significant reality gap: the rendered outputs lack photorealism and often fail to represent the visual complexity of real-world driving scenes. 
Reconstruction-based simulators, such as PVG~\cite{chen2023periodic}, OmniRe~\cite{chen2024omnire}, and S3Gaussian~\cite{huang2024s3gaussian} reconstruct driving scenes from recorded trajectories. 
However, they can only render high-quality views along pre-recorded trajectories, limiting their ability to support free-viewpoint simulation.
FreeSim~\cite{fan2024freesim}, DriveDreamer4D~\cite{zhao2024drive}, and ReconDreamer~\cite{Ni2024ReconDreamerCW} use diffusion-based generative models to synthesize novel views and reconstruct 4D dynamic scenes from both real and synthesized images.
However, these methods often struggle to maintain view consistency or preserve details when combining generated and real views, resulting in suboptimal reconstructions and blurred visual results. Additionally, they generally lack the capability for scene editing.
% While promising, these methods face significant challenges in view consistency, especially in details: the details of the same regions may be generated differently across trajectories, leading to blurring when reconstructing the scenes using generated and real views. Furthermore, these systems typically lack scene editing capabilities, making it difficult to manipulate environmental conditions after reconstruction.
% To overcome the off-trajectory limitation, recent approaches like FreeSim~\cite{fan2024freesim}, DriveDreamer4D~\cite{zhao2024drive}, and ReconDreamer~\cite{Ni2024ReconDreamerCW} use diffusion-based generative models to synthesize novel views and reconstruct 4D dynamic scenes from both real and synthesized images. While promising, these methods face significant challenges in view consistency, especially in details: the details of the same regions may be generated differently across trajectories, leading to blurring when reconstructing the scenes using generated and real views. Furthermore, these systems typically lack scene editing capabilities, making it difficult to manipulate environmental conditions after reconstruction.
Recent world models, such as DriveDreamer-2~\cite{zhao2025drivedreamer} and MaskGWM~\cite{ni2025maskgwm}, attempt to model dynamic scenes by generating videos. 
However, unlike reconstruction-based methods, they can only control the camera pose in a coarse manner, and 3D consistency cannot be guaranteed. As a result, they are unable to support accurate free-viewpoint simulation.

% what about the video diffusion based world model?

% To end this, we propose GA-Drive, adopting Geometry-Appearance Decoupling and Diffusion-Based Rendering that can simulate photorealistic off-trajectory views and achieve appearance editing through text prompts. First, we modify the OmniRe [xxx], a 4D Gaussian Scene Graph representation to have a accurate geometry reconstruction. Since the original OmniRe only use LiDAR-based depth supervision, it is constrained by the limited sensing range and the sparsity of the LiDAR points. Consequently, it provides insufficient depth information for distant regions and the upper parts of the driving scene. To solve this, we add an aligned depth maps to supervise the depth maps rendered from the OmniRe. After having the clear 4D geometry, given a target view, we can project each pixel of the target view into the space according to the depth of the target view to form a point cloud. Then, the point cloud can be projected onto the known images to sample the color. Since some regions appearing in the known image can be blocked I the target view. Thus, a visibility check is needed to mask these blocked regions. Then, we can have the novel pseudo-views. We use a video diffusion model to render the novel trajectory views by inputting these novel pseudo-views. To train the video diffusion model, we develop a data augmentation pipeline that can train the diffusion model without training a 4D representation unlike, dd4d, recondreamer, and freesim. The reason of doing so is the following: 1. 4D reconstruction is time-consuming, 1.5 hours for a scene. 2. 4D reconstruction highly relies on accurate camera pose and lidar data. Thus, the dataset can not be scaled up. 3. Since we want to edit the scene, including weather editing, the current driving datasets nuscens, waymo,[xxx].. contain very few data on the extreme weather like foggy, snowy, etc. To achieve editing through text prompts, an image-to-image method like InstructPix2Pix can be used as an appearance editor. After altering the appearance of the original recorded image, the appearance of the novel views will also be changed accordingly. Because we adopt the Geometry-Appearance Decoupling, although the editing is done on the image separately, the underlying geometry is still consistent across frames, and we utilize video diffusion model not image diffusion model, which can mitigate the subtle detailed difference and coherence between frames. Furthermore, we propose progressive trajectory expansion strategy to solve limitations: the video diffusion model can only generate a limited length of video. Large trajectory shift have a large blocked region which may causes difficulty for the video diffusion model to inpaint the large blocked region still with high quality. The progressive trajectory expansion strategy contains 2 parts to solve the above problem. 1. In the data augmentation pipeline, we do not add a block mask to the first frame, which will let the diffusion model generate the block region according to the first frame. Then, when we generated the second video segment, we can use the last frame of the previous video segment as the first pseudo image frame. In this way, the two video segments can be concatenated seamlessly. 2. Instead of directly shifting to the target trajectory with a large deviation, we progressively shift the trajectory toward the target trajectory step by step. After each step, we generate the pseudo novel trajectory and use the diffusion model to generate the photorealistic trajectory. Then, this novel trajectory will be recorded as a known trajectory. And the pseudo image of the next step will take reference to this trajectory as well as other known trajectories. In such way, the large trajectory shifts can be achieved. 

% We begin by improving upon a reconstruction-based OmniRe~\cite {chen2024omnire} to achieve more accurate geometry reconstruction by introducing dense depth supervision and depth distortion regularization. 
% Dense depth supervision enhances the reconstructed geometry, particularly in areas where LiDAR data is sparse or unavailable. Depth distortion regularization enforces the floating gaussions to cluster with other gaussions to form a clear surface. 

To this end, we propose GA‑Drive, a novel photorealistic driving simulator that supports free‑viewpoint simulation and appearance editing via geometry–appearance decoupling and diffusion‑based generation.
Our key insight is to decouple geometry from appearance. In driving scenes with limited viewpoints, RGB images cannot uniquely determine 3D geometry. Fitting appearance sacrifices geometric accuracy by introducing floaters or distortions, which cause obvious artifacts on novel trajectories. By separating these components, we enable the geometry module to focus on accurate 3D structure while delegating appearance synthesis to a diffusion model.


% Our geometry module is built on a 4D reconstruction method OmniRe~\cite{chen2024omnire}. We introduce dense depth supervision and depth‑distortion regularization to obtain a more accurate 4D geometry reconstruction. 
% % 为了告诉这个diffusion足够的信息去生成3d consistency的novel view, 我们利用我们准确的几何重建去生成了一个novel pseudo-view,这个做一个diffusion condition。
% Given this clean reconstruction, we synthesize novel pseudo-views by casting rays from the novel pose into 3D space to form a point cloud, then projecting the points onto captured views to retrieve their colors. 
% %To handle occlusions, we perform a multi‑view visibility check and mask pixels not visible in any source view.
% Although these novel pseudo-views may contain holes or masked regions, they serve as a strong prior that our video diffusion model leverages to generate photorealistic, 3D‑consistent novel views.

Our geometry module builds on OmniRe~\cite{chen2024omnire}, a 4D reconstruction method. We introduce dense depth supervision and depth-distortion regularization to obtain more accurate geometry. To provide the diffusion model with sufficient geometric and appearance cues for generating 3D-consistent novel views, we leverage this accurate reconstruction to synthesize novel pseudo-views that encode correct 3D spatial relationships and occlusion patterns. Specifically, we cast rays from the target novel pose into the reconstructed 3D space, forming a point cloud, then project these points onto captured source views to retrieve their colors. Although the resulting pseudo-views may contain holes or occluded regions, they serve as a strong prior that guides our video diffusion model to generate photorealistic, 3D-consistent novel views.


% To enable robust video editing under diverse weather conditions, it is critical to train the video diffusion model on a wide range of environmental scenarios. However, existing driving video datasets such as Waymo~\cite{Sun_2020_CVPR}, nuScenes~\cite{caesar2020nuscenes}, and others contain only a limited number of samples captured under extreme weather conditions, such as snow or fog.

% Such reconstruction-based training data generation requires high-quality multiview posed data and has limited scalability.
% As a result, web-scale diverse data can not be leveraged,
% % This approach is adopted in methods like FreeSim~\cite{fan2024freesim}. However, the effectiveness of such training is inherently constrained by the diversity of the underlying dataset. 
% % Since extreme-weather scenes are underrepresented, the resulting pseudo-images also lack such diversity, 
% limiting the model’s generalizability.
% A common strategy for generating training pairs to train the diffusion model involves 4D reconstruction, where a subset of images is used to reconstruct the 4D scene, and the remaining images serve as ground truth for the pseudo views rendered from the partial reconstruction, like Freesim and recondeamer. 
% To train the video diffusion model, we introduce an image augmentation pipeline that simulates pseudo-view patterns, eliminating the need for expensive 4D reconstruction to generate training pairs as in Freesim~\cite{fan2024freesim}. 

% A common strategy for generating training pairs to train the diffusion model involves 4D reconstruction as in Freesim~\cite{fan2024freesim}. However, 4D reconstruction is time-consuming, thus it is difficult to get large pseudo novel images and corresponding ground truth images by aformentioned method. To solve it, we introduce an image augmentation pipeline that simulates pseudo-view patterns, eliminating the need for expensive 4D reconstruction.


% A common strategy for generating training pairs for novel trajectory generation involves 4D reconstruction, as in FreeSim~\cite{fan2024freesim}. 
% However, obtaining a large number of reconstructed scenes is extremely time-consuming, making it difficult to scale up the training data.


% The training of our video diffusion model requires pseudo-views as conditioning frames and corresponding ground-truth images as supervision targets. However, public driving datasets typically offer only a single trajectory per scene, since the vehicle cannot drive along multiple paths simultaneously. Consequently, while the captured views from that trajectory can serve as ground-truth targets, there is no second trajectory to serve as sources for synthesizing pseudo-views of the recorded path.
% To address this difficulty, we introduce a pseudo-view simulation pipeline that generates pseudo-views from the single available trajectory, circumventing the need for multiple simultaneously recorded trajectories. Simulated pseudo-views serve as conditioning frames for the video diffusion model, while the original captured views serve as ground-truth supervision. Moreover, this approach enables scalable training on large-scale datasets such as Waymo~\cite{Sun_2020_CVPR} and OpenDV~\cite{gao2024vista}, equipping the model to handle diverse driving scenarios, including challenging conditions like fog and snow.


% However, training our video diffusion model presents a key challenge: it requires pseudo-views as conditioning frames and corresponding ground-truth images as supervision targets. Ideally, with multiple trajectories in the same scene, one could use captured views from trajectory A to synthesize pseudo-views of trajectory B (via geometric projection), then supervise the model using real captured views from trajectory B. Unfortunately, public driving datasets typically offer only a single trajectory per scene, since a vehicle cannot traverse multiple paths simultaneously. This means that while the captured views can serve as ground-truth targets, there is no second trajectory to provide source views for synthesizing the necessary pseudo-view conditions.

However, training our video diffusion model presents a key challenge. Ideally, with multiple trajectories in the same scene, we could use views from trajectory A to synthesize pseudo-views of trajectory B, then supervise the model with real views from trajectory B. Unfortunately, public driving datasets offer only a single trajectory per scene. While captured views can serve as ground-truth supervision, there is no second trajectory to provide sources for synthesizing pseudo-view conditions.

To address this challenge, we introduce a pseudo-view simulation pipeline that simulates training pairs from the single available trajectory, circumventing the need for multiple simultaneously recorded trajectories. Simulated pseudo-views serve as conditioning frames for the video diffusion model, while the original captured views serve as ground-truth supervision. Moreover, this approach enables scalable training on large-scale datasets such as Waymo~\cite{Sun_2020_CVPR} and OpenDV~\cite{gao2024vista}, equipping the model to handle diverse driving scenarios, including challenging conditions like fog and snow.


% To enable long video generation, we train our model to generate video segments conditioned on pseudo-novel views, using the last frame of each segment as the starting frame for the next. This segment-wise approach ensures temporal consistency across the generated novel-view video.
% Since video diffusion models typically have limited generation lengths, we train our model to generate video segments conditioned on novel pseudo-views, with a particular emphasis on the first frame. Specifically, the first frame of the pseudo views is set as the last frame of the previous generated segment, ensuring temporal consistency across consecutive segments. In this way, our video diffusion model can generate long, novel-view videos in a segment-wise manner.

% Since all geometric information is encapsulated in the 4D geometry representation and the recorded views serve solely as a source for sampling color or appearance, geometry and appearance are effectively decoupled.
Since all geometric information is encapsulated in the 4D geometry representation while recorded views serve solely as appearance sources, geometry and appearance are effectively decoupled.
% In previous reconstruction-based methods, geometry and appearance are coupled and usually conflict with each other to some degree, which is caused by sensor noise and distortion.
% After decoupling, our method can produce geometry-consistent, photorealistic appearances for novel trajectory views. This decoupling also enables flexible appearance editing without affecting the underlying geometry. Once the appearance of the original recorded videos is modified, the appearance of the generated novel views is updated accordingly. The editing of the original videos can be achieved using state-of-the-art video editing methods, such as VACE~\cite{vace}, Instruct-V2V~\cite{cheng2023consistent}, or other advanced techniques.
After decoupling, our method produces geometry-consistent, photorealistic appearances for novel trajectory views. This decoupling also enables flexible appearance editing without affecting the underlying geometry. Once the appearance of the original recorded videos is modified, the appearance of all generated novel views across arbitrary trajectories is updated accordingly, ensuring global consistency. The editing of the original videos can be achieved using state-of-the-art video editing methods, such as VACE~\cite{vace}, Instruct-V2V~\cite{cheng2023consistent}, or other advanced techniques.




% For instance, Instruct-V2V can modify the weather based on a text prompt, while VACE can change the appearance of a specific object.

% Thus,50we model support text-driven appearance editing by editing the original video. Once the original51recorded videos are edited according to a user-defined prompt, the appearance of the generated novel52views changes accordingly.53
% To train the video diffusion model, we propose a video augmentation pipeline that generates training data without expensive reconstruction, scaling up the training data. The video augmentation pipeline consists of masking regions according frame's depth maps to mimic the occluded regions and other commonly used augmentation method like add gaussion noise and random maskes. 

% We collect the Waymo~\cite{Sun_2020_CVPR} and Opendv ~\cite{gao2024vista} dataset as the training data. 

% After training with diverse data, it becomes capable of processing pseudo views across a wide range of driving scenarios, including extreme weather conditions such as fog and snow. Video diffusion models typically have limited generation lengths. To address these limitations, we 


% we propose a Trajectory Expansion Strategy, comprising two modules: the Head-to-Tail Trajectory Expansion module and the Progressive Trajectory Shifting module. Together, these approaches enable the generation of arbitrarily long videos with substantial trajectory shifts.

% Thus, our model support text-driven appearance editing by editing the original video.

% % We utilize  InstructPix2Pix~\cite{brooks2022instructpix2pix}, to edit images. 
% Once the original recorded videos are edited according to a user-defined prompt, the appearance of the generated novel views changes accordingly. Thanks to our video diffusion model trained with diverse data, we can achieve appearance editing, such as altering the weather to foggy or snowy conditions.
% % simulates the characteristic patterns observed in pseudo-trajectory images, enabling the training of video diffusion models using only videos, without requiring 4D reconstruction.


% Unlike previous approaches such as DD4D, ReconDreamer, and FreeSim, which require joint training with a 4D reconstruction module, our approach decouples geometry learning from appearance synthesis, allowing the video diffusion model to be trained independently via a dedicated data augmentation pipeline. This design choice offers several benefits. First, 4D reconstruction is computationally intensive, often taking over an hour per scene, which limits scalability. Second, it requires accurate camera poses and dense LiDAR scans, which are not always available or reliable. Third, because we aim to support appearance editing—especially in rare or extreme weather conditions such as fog or snow—existing driving datasets like nuScenes, Waymo, and others [XXX] offer limited data diversity, particularly in these scenarios. Our method bypasses these limitations by enabling scalable training and generation without dependence on dense 4D reconstructions or perfectly aligned sensor inputs.

% Because we adopt Geometry-Appearance Decoupling, the underlying 4D geometry remains the same across frames, ensuring spatial consistency even as appearance varies. 
% Moreover, since the editing is handled within a video diffusion model, rather than frame-by-frame with an image diffusion model, temporal coherence is preserved and subtle inconsistencies across frames are mitigated.
% A remaining challenge is that video diffusion models typically have limited generation lengths and may struggle with large trajectory shifts, especially when these shifts introduce significant occlusions. To address these limitations, we propose a Trajectory Expansion Strategy, comprising two modules: the Head-to-Tail Trajectory Expansion module and the Progressive Trajectory Shifting module. Together, these approaches enable the generation of arbitrarily long videos with substantial trajectory shifts.

% In this strategy, we can generate infinite-length videos with the Head-to-Tail Trajectory Expansion strategy

% the first frame of each generated video segment is left unmasked, allowing the diffusion model to ground the content generation based on the initial visible context. When generating the next segment, we use the last frame of the previous segment as the starting pseudo frame, enabling seamless transitions across segments. Additionally, instead of jumping directly to distant target views, we incrementally shift the camera trajectory toward the desired path in small steps. After each step, we generate pseudo views and refine them using the diffusion model. The new trajectory segment is then added to the known view set, so that subsequent steps benefit from increasingly rich visual context. Through this iterative process, GA-Drive can handle large viewpoint changes while maintaining high visual quality and consistency.

We summarize our contributions as follows
\begin{enumerate}
    \item \textbf{Geometry-Appearance Decoupling}: A novel framework that separates geometry and appearance for diffusion-based generation in driving simulators. This feature enables flexible appearance editing.
    \item \textbf{Pseudo-view Simulation Pipeline}: A pipeline that generates training pairs for video diffusion models without requiring datasets with multiple simultaneously recorded trajectories.
    % \item \textbf{Segment-wise Video Generation Framework}: that generates long, temporally consistent novel-view videos in a segment-wise manner.
    \item The proposed framework substantially outperforms existing methods for the novel trajectory synthesis task in terms of NTA-IoU, NTL-IoU, and FID scores.

\end{enumerate}